\documentclass[border=3pt]{standalone}
\usepackage{tikz}
\usetikzlibrary{positioning,fit,arrows.meta}
\usepackage{xeCJK}
\setCJKmainfont{KaiTi_GB2312}
\setCJKsansfont{SimHei}
\begin{document}
\begin{tikzpicture}[font=\small,
  layer/.style={draw,thick,rounded corners=2pt,align=left,minimum height=1.0cm,text width=12.4cm,fill=blue!5},
  ar/.style={-{Stealth[length=2mm]},thick}]
  \node[layer] (l1) {\textbf{(1) 冻结计划}：主类对、样本量、停止条件、Holm 家族在采集前写定};
  \node[layer,below=0.34cm of l1] (l2) {\textbf{(2) 被测模块}：M1--M5，登记源摘要、产物摘要、编译器与选项};
  \node[layer,below=0.34cm of l2] (l3) {\textbf{(3) 测量层}：单次解封装计时循环、锁频控制与恢复核对、三次独立运行、原始值落盘};
  \node[layer,below=0.34cm of l3,minimum height=2.0cm] (l4) {};
  \node[anchor=north west,font=\small\bfseries] at ([xshift=0.1cm,yshift=-0.08cm]l4.north west) {(4) 分析层：三路独立登记};
  \node[draw,fill=white,rounded corners=2pt,text width=3.6cm,align=center,font=\footnotesize] at ([xshift=-4.15cm,yshift=-0.32cm]l4.center) {统计判定\\ Welch / Holm / 合并 / MDE};
  \node[draw,fill=white,rounded corners=2pt,text width=3.6cm,align=center,font=\footnotesize] at ([yshift=-0.32cm]l4.center) {定位工具\\ 指令、分支与访存};
  \node[draw,fill=white,rounded corners=2pt,text width=3.6cm,align=center,font=\footnotesize] at ([xshift=4.15cm,yshift=-0.32cm]l4.center) {族外交叉检验\\ RTLF / dudect / Mona};
  \node[layer,below=0.34cm of l4,fill=green!8] (l5) {\textbf{(5) 结果登记层}：检出 / 未检出（附 MDE）/ 未执行 / 不适用，带运行编号与证据编号};
  \foreach \a/\b in {l1/l2,l2/l3,l3/l4,l4/l5} \draw[ar] (\a) -- (\b);
  \node[draw,dashed,rounded corners=4pt,fit=(l1)(l5),inner sep=0.30cm] {};
  \node[font=\scriptsize,align=center,text width=12.4cm,below=0.52cm of l5,fill=white,rounded corners=2pt] (proc)
    {同一套规程贯穿五层：同一产物、同一输入类、同一停止条件；结论只在这套规程内成立};
  \node[draw,dashed,rounded corners=2pt,font=\footnotesize,align=center,text width=12.4cm,below=0.35cm of proc,fill=yellow!8]
    {四层证据：机制层（定位）→ 热点层（函数级计时）→ 端到端层（完整解封装）→ 防护复测层（同规程重测）};
\end{tikzpicture}
\end{document}
